\begin{table}[h!]
    \begin{tabular}{| p{0.20\textwidth} | p{0.35\textwidth} | p{0.35\textwidth} |}
    \hline
    \textbf{Pendekatan} 
    & \textbf{Kelebihan} 
    & \textbf{Keterbatasan} \\
    \hline

    \textbf{Deep learning / Transformer} 
    & Representasi semantik kuat, hasil unggul di \textit{dataset} besar	 
    & Membutuhkan GPU dan \textit{dataset} besar; tidak sesuai batasan TA \\
    \hline

    \textbf{ML klasik berbasis konten (TF–IDF)} 
    & Efisien, stabil, banyak dipakai dalam penelitian hoaks Indonesia	 
    & Terbatas pada teks; tidak menangkap konteks sumber \\
    \hline

    \textbf{ML klasik + metadata (mis. source)} 
    & Menambah sinyal kredibilitas media; ruang lingkup tetap terbatas
    & Kualitas metadata menentukan efektivitas pendekatan \\
    \hline
    \end{tabular}
\caption{Perbandingan pendekatan untuk deteksi hoaks}
\label{tbl:perbandingan-pendekatan}
\end{table}



